\begin{textsample}{POS Dim 3 – rs – Score 12.00 – rs/rs\_inf\_19.txt}  \label{ex:f3_pos_230}
13.
CONSIDERAÇÕES FINAIS A construção de um currículo que considerasse as especificidades do nosso Estado é um \textbf{desejo} e uma reivindicação dos educadores gaúchos.
Nesse documento de referência curricular, foi possível considerar os aspectos locais e regionais nos Objetivos de aprendizagem e desenvolvimento e a valorização da diversidade cultural do território, a partir das contribuições dos educadores de todo o Estado, via plataforma digital.
Este documento curricular representa um marco histórico de reconhecimento da Educação \textbf{Infantil} como primeira etapa da Educação Básica, com identidade própria e com um currículo que afirma os direitos e as especificidades do desenvolvimento de \textbf{bebês}, \textbf{crianças} bem \textbf{pequenas} e \textbf{crianças} \textbf{pequenas}.
Ao finalizar esse documento, é importante destacar a necessidade da consolidação, em cada instituição, de um currículo brincante, relacional, capaz de produzir significado e \textbf{acolher} a ludicidade, a criatividade, a autoria e a \textbf{curiosidade} de \textbf{crianças} e \textbf{adultos}.
Um currículo que coloque no centro do processo educativo as \textbf{crianças}, as interações e a \textbf{brincadeira}, e que promova possibilidades de desenvolvimento e aprendizagem por meio de múltiplas linguagens.
Um currículo fundamentado na garantia dos direitos das \textbf{crianças}, articulados aos Campos de Experiências, organização curricular que parte da concepção de uma pedagogia relacional, em que o conhecimento é produzido na interação entre as \textbf{crianças} e o mundo, entre os \textbf{adultos} e as \textbf{crianças}, entre as \textbf{crianças} e outras \textbf{crianças}, aberto para a complexidade que é conhecer -se e conhecer o mundo.
Um currículo que busca dar sentido à variedade de experiências que as \textbf{crianças} experimentam cotidianamente na escola.
A complexidade de um currículo fortemente marcado pela intencionalidade educativa, pela presença atenta, sensível e interessada de um professor que, junto às \textbf{crianças}, se lança a experimentar e descobrir o mundo.
Nesse sentido, o grande desafio é o de promover uma educação que integre, nas práticas cotidianas, a atenção aos \textbf{cuidados}, à saúde, à proteção, assim como as aprendizagens dos repertórios culturais e da diversidade de linguagens a que as \textbf{crianças} têm direito.
Trata -se de comprometer -se com escolhas discutidas, refletidas e definidas coletivamente, entre os educadores, em parceria com as famílias e a partir da \textbf{escuta} atenta das \textbf{crianças}, para fundamentar a construção de um projeto educacional planejado e desenvolvido pelos envolvidos no objetivo comum que é oferecer educação de qualidade aos \textbf{bebês}, às \textbf{crianças} bem \textbf{pequenas} e às \textbf{crianças} \textbf{pequenas}, ao encontro de aprendizagens significativas na Educação \textbf{Infantil}.

% matched lemmas: acolher, adulto, bebê, brincadeira, criança, cuidado, curiosidade, desejo, escuta, infantil, pequeno
\end{textsample}
